\subsection{Simplifying the $t$-\toughpair}\label{sec:firstramsey}

In this section we prove Lemma~\ref{lem:pre-cleaning} restated below.
\firstramsey*

\begin{proof}
Let $(E_1,E_2)$ be an $ \mathcal{R}(2t,9s)$-\toughpair\ in $D$.
Let $t' =\mathcal{R}(2t,9)$.
 Fix an arbitrary ordering of the edges of $E_1$ and $E_2$.  Let $E_1=(e^1_1, \ldots, e^1_{t'})$ and let $E_2=(e^2_1, \ldots, e^2_{t'})$.
Fix $i \in [2]$. We will use Ramsey arguments to prove the lemma. Towards this, construct an auxiliary undirected, complete, edge-colored graph $\aux{i}$ as follows. The vertex set of $\aux{i}$ corresponds to the edges of $E_i$, that is, for each $e \in E_i$ there is a vertex corresponding to $e$.
For the sake of simplicity we denote the vertex of $\aux{i}$ that correspond to the edge $e$ of $E_i$, by $e$ itself.
The coloring function $\col_i$ on the edges of $\aux{i}$ (equivalently on the pair of distinct edges of $E_i$) is defined based on the following. Fix any two distinct edges $e=(u,v),e'=(u',v') \in E_i$. The coloring function $\col_i(e,e')$ is defined based on the existence of the edges $(u,v')$ and $(u',v)$ in $D$ and $D^{\star}$. Below we describe $\col_i(e,e')$.

\begin{enumerate}
\item If $(u,v') \in E(D)$, and
\begin{enumerate}
\item $(u',v) \in E(D)$, then $\col_i(e,e') =(1,1)$,
\item $(u',v) \in E(D^{\star}) \setminus E(D)$, then $\col_i(e,e') =(1,2)$,
\item $(u',v) \not \in E(D)$, $(u',v) \not \in E(D^{\star})$, then $\col_i(e,e') = (1,3)$.
\end{enumerate}

\item If $(u,v') \in E(D^{\star}) \setminus E(D)$, and
\begin{enumerate}
\item $(u',v) \in E(D)$, then $\col_i(e,e') =(2,1)$,
\item $(u',v) \in E(D^{\star}) \setminus E(D)$, then $\col_i(e,e') =(2,2)$,
\item $(u',v) \not \in E(D)$, $(u',v) \not \in E(D^{\star})$, then $\col_i(e,e') = (2,3)$.
\end{enumerate}

\item If $(u,v') \not \in E(D)$, $(u,v') \not \in E(D^{\star})$, and
\begin{enumerate}
\item $(u',v) \in E(D)$, then $\col_i(e,e') =(3,1)$,
\item $(u',v) \in E(D^{\star}) \setminus E(D)$, then $\col_i(e,e') =(3,2)$,
\item $(u',v) \not \in E(D)$, $(u',v) \not \in E(D^{\star})$, then $\col_i(e,e') = (3,3)$.
\end{enumerate}
\end{enumerate}

We now use a Ramsey argument on the $t'$-vertex graph $\aux{i}$ with edge-coloring function $\col_i$ (that uses at most $9$ different colors). Using Ramsey arguments, we conclude that there exists a monochromatic clique in $\aux{i}$ of size $2t$. Below we show how one can get one of the three outcomes in the lemma statement based of the color of the monochromatic clique. Say edges of $E_i$ that correspond to this monochromatic clique in $\aux{i}$ are 
$E'_i=\{e^i_{j_1}, e^i_{j_2}, \ldots , e^i_{j_{t}}\}$ such that $j_1 < \ldots < j_{t}$ (recall we fixed an ordering of the edges of $E_i$). For the ease of notation later, we assume that the sets $E'_i$ are ordered with the ordering as described in the previous line.

\begin{claim}\label{claim:ramsey-orderedpair}
If for each $i \in [2]$, the color of the monochromatic clique in $\aux{i}$ is $(j,3)$ or $(3,j)$ for any $j \in [3]$, then $(E'_1,E'_2)$ form an ordered $2t$-\toughpair\ in $D$.
\end{claim}

\begin{claimproof}
Since $(E_1,E_2)$ is a \toughpair\ and $E'_i \subseteq E_i$, $(E'_1,E'_2)$ is also a \toughpair\ (of size $t$). If $\col_i$ colors the clique corresponding to the edges of $E'_i$ with the color $(j,3)$, then consider the ordering $(e^i_{j_1}, e^i_{j_2}, \ldots , e^i_{j_{t}})$ of $E'_i$, otherwise (when $\col_i$ colors the monochromatic clique with color $(3,j)$), then consider the ordering $(e^i_{j_{t}}, \ldots , e^i_{j_1})$ of $E'_i$. Then from the description of the definition of $\col_i$ corresponding to the case when $\col_i$ takes value $(3,j)$ or $(j,3)$, one concludes that $(E'_1,E'_2)$ is indeed an ordered $t$-\toughpair\ with the ordering described above.
\end{claimproof}
Note that if the graph has an ordered $2t$-\toughpair, then it also has an ordered $t$-\toughpair.

\begin{observation}\label{obs:ramsey-biclique}
It is easy to observe that if the color of the monochromatic clique in $\aux{i}$ is $(1,1)$, then $(\tail(E'_i),\head(E'_i))$ form a $2t$-biclique in $D$. If the color is $(2,2)$, then $D$ contains the $(\tail(E'_i),\head(E'_i))$-induced matching and $D^{\star}$ contains the $(\tail(E'_i),\head(E'_i))$-biclique (thus, the third outcome of the lemma holds). 
\end{observation}

Note that if the graph has $2t$-biclique, then it also has an ordered $t$-biclique.

\begin{claim}\label{claim:ramsey-induced-skew}
If there exists $i \in [2]$ such that the color of the monochromatic clique in $\aux{i}$ is $(1,2)$ or $(2,1)$, then $H$ contains a $t$-biclique.
\end{claim}

\begin{claimproof}
Suppose that the color of the monochromatic clique in $\aux{i}$ is $(1,2)$. Recall that $E'_i =((e^i_{j_1}, e^i_{j_2}, \ldots , e^i_{j_{t}}))$ is the ordered set of edges that correspond to the monochromatic clique. The $(\tail(\{e^i_{j_1}, \ldots , \allowbreak e^i_{j_{t}}\}), \allowbreak \head(\{e^i_{j_{t +1 }}, \ldots, \allowbreak e^i_{j_{2t}}\}))$ form a $t$-biclique in $D$.

Similarly, if the color of the monochromatic clique in $\aux{i}$ is $(2,1)$, then $(\tail(\{e^i_{j_{t+1}}, \ldots, \allowbreak e^i_{j_{2t}}\}), \allowbreak  \head(\{e^i_{j_1}, \ldots, \allowbreak e^i_{j_{t}}\}))$ form a $t$-biclique in $H$.
\end{claimproof}

Thus, from Claims~\ref{claim:ramsey-induced-skew} and~\ref{claim:ramsey-orderedpair}, and Observation~\ref{obs:ramsey-biclique}, the lemma follows.
\end{proof}